\section{模型假设与符号说明}

\subsection{模型假设}

在题目给定的设备参数和交易规则基础上，为统一四个问题中的电量流向、时间尺度和信息使用范围，作如下假设。

\begin{enumerate}
    \item \textbf{光伏优先直接满足小区负载。}
    当光伏功率超过当前负载时，富余部分可用于储能；外网购电既可直接供给负载，也可用于给储能设备充电。题目未给出售电规则，因此未被负载或储能利用的光伏不计售电收益。

    \item \textbf{以10分钟为一个调度时段。}
    认为每个采样点的负载功率、光伏功率和电价在其对应的10分钟区间内保持不变，并据此将功率转换为区间电量。该离散化与题目要求的10分钟购电结果保持一致。

    \item \textbf{储电状态按时间连续衔接。}
    问题一按照题意要求日初和日末储电量相等；问题二至问题四则将前一天24:00的实际剩余电量作为下一天0:00的初始电量，不再每日重置储能状态。

    \item \textbf{决策只使用当时已经获得的信息。}
    问题二制定零点计划时，只使用此前已经发生的负载、光伏和储电状态；问题三只使用对应时刻已经发布的光伏预报；问题四中未来实际电价同样不能提前使用。由此避免在模型中引入未来信息。

    \item \textbf{订单调整以零点原计划为统一基准。}
    在问题三及问题四的日内调整分支中，已经执行的订单不再修改；同一时段经过多次调整时，只计算最终订单相对零点原计划的净变化，避免重复计算调整费用。

    \item \textbf{不额外考虑题目未给出的储能损耗因素。}
    除题目规定的充、放电效率外，不再引入电池自放电、寿命折旧等额外成本。这样既避免人为增加缺乏参数支持的费用项，也使目标函数集中于题目要求的购电成本。

\end{enumerate}


\subsection{符号说明}

设 $d$ 表示日期，$k$ 表示一天中的第 $k$ 个10分钟时段，其中

\begin{equation}
k=1,2,\ldots,144,
\qquad
\Delta t=\frac{1}{6}\ \mathrm{h}.
\end{equation}

$S_{d,0}$ 表示第 $d$ 天0:00的储电量，$S_{d,k}$ 表示第 $k$ 个时段结束后的储电量。问题一仅研究一个代表日，因此相应公式中可省略日期下标 $d$。

主要公共符号如表所示。各问中特有的预测误差、准备量、订单调整量及价格模型参数，在对应模型中首次出现时再单独定义。

\begin{table}[H]
    \centering
    \caption{主要符号及其含义}
    \begin{tabular}{ccl}
        \hline
        符号 & 含义 & 单位或取值 \\
        \hline
        $d$ & 日期下标 & -- \\
        $k$ & 日内10分钟时段下标 & $1,\ldots,144$ \\
        $\Delta t$ & 单个调度时段长度 & $1/6\,\mathrm{h}$ \\
        $L_{d,k}$ & 小区负载功率 & kW \\
        $P^{PV}_{d,k}$ & 光伏发电功率 & kW \\
        $E^L_{d,k}$ & 第 $k$ 时段负载电量 & kWh \\
        $E^{PV}_{d,k}$ & 第 $k$ 时段光伏电量 & kWh \\
        $p_{d,k}$ & 对应时段外网电价 & 元/kWh \\
        $m_{d,k}$ & 光伏直接供给负载的电量 & kWh \\
        $n_{d,k}$ & 光伏直供后的负载缺口 & kWh \\
        $r_{d,k}$ & 光伏直供后的富余电量 & kWh \\
        $q_{d,k}$ & 普通外网购电量 & kWh \\
        $x_{d,k}$ & 外网购电中用于储能充电的电量 & kWh \\
        $y_{d,k}$ & 富余光伏用于储能充电的电量 & kWh \\
        $z_{d,k}$ & 储能设备向负载输出的电量 & kWh \\
        $S_{d,k}$ & 第 $k$ 时段结束后的内部储电量 & kWh \\
        $w_{d,k}$ & 未被负载或储能利用的光伏电量 & kWh \\
        $S_{\min}$ & 储电量运行下限 & 1200 kWh \\
        $S_{\max}$ & 储电量运行上限 & 10800 kWh \\
        $P_{\max}$ & 最大充、放电功率 & 5000 kW \\
        $\eta_c$ & 充电效率 & 0.90 \\
        $\eta_d$ & 放电效率 & 0.90 \\
        \hline
    \end{tabular}
\end{table}

上述符号构成四个问题共用的电量流和储能状态基础，其中设备容量、功率和效率均直接采用题目给定参数，而非额外假设。